\newcommand{\C}{\mathbb{C}}
\newcommand{\R}{\mathbb{R}}
\newcommand{\Z}{\mathbb{Z}}
\newcommand{\Tr}{\operatorname{Tr}}
\newcommand{\Impart}{\operatorname{Im}}
\newcommand{\AF}{\operatorname{AF}}
\newcommand{\AX}{\operatorname{AX}}
\newcommand{\GL}{\operatorname{GL}}

\examyear{2012}
\examera{平成24年}
\examfield{幾何}
\examgroup{B}
\examnumber{5}
\examtitle{平成24年 B 第5問}
\examtag{リーマン幾何}
\examtag{体積要素}
\examtag{ベクトル場}

開円板 \(D=\{(x,y)\in\R^2\mid x^2+y^2<1\}\) にリーマン計量
\[
  ds^2=\frac{4(dx^2+dy^2)}{(1-x^2-y^2)^2}
\]
を入れる。

\begin{enumerate}
  \item このリーマン計量についての体積要素 \(\omega\) を求めよ。
  \item \(0<s<1\) に対する \(C^\infty\) 関数 \(f(s)\) であって次の条件を満たすものをすべて求めよ：
  \(D\) から原点 \(O\) を除いた領域 \(D_0\) においてベクトル場 \(X\) を
  \[
    X=f(x^2+y^2)\left(x\frac{\partial}{\partial x}+y\frac{\partial}{\partial y}\right)
  \]
  によって定める。このとき \(t\ge0\) をパラメータとする \(C^\infty\) 写像の族
  \(\varphi_t:D_0\to D_0\) が存在し，次を満たす。
  \[
  \begin{cases}
    \varphi_0(p)=p & (p\in D_0),\\
    \dfrac{d}{dt}\varphi_t(p)=X_{\varphi_t(p)} & (p\in D_0,\ t\ge0),\\
    \varphi_t^*\omega=\omega & (t\ge0).
  \end{cases}
  \]
\end{enumerate}